\section{The FFN as Belief Update}
\label{sec:ffn-or}

Once attention has assembled the relevant evidence into the residual
stream, the feed-forward network performs the update step. In the exact
constructive setting, this update is the binary belief-propagation rule
for combining independent evidence sources.

The central formula is:
\[
\updateBelief(m_0, m_1)
  = \frac{m_0 m_1}{m_0 m_1 + (1-m_0)(1-m_1)}
  = \sigma\!\left(\logit(m_0) + \logit(m_1)\right).
\]
With sigmoid activation and BP weights ($w_0 = w_1 = 1$, $b = 0$), the
FFN computes this exactly from the gathered beliefs in dimensions~4
and~5, writing the result to dimension~0. This is the Bayesian update
rule for two independent binary evidence sources in the constructive
model~\cite{coppola2026transformers}.

\subsection*{The Boolean Structure Inside the Update}

The formula contains the same structure that appears in classical
Boolean reasoning, but in probabilistic form. The term $m_0 m_1$ is the
joint probability that both independent inputs are true. In log-odds
space, this multiplication becomes addition:
\[
\logit(\updateBelief(m_0, m_1)) = \logit(m_0) + \logit(m_1).
\]

Independent evidence adds. The sigmoid converts the log-odds sum back to
probability space. The FFN is therefore a weight-of-evidence combiner in
the sense of Turing and Good~\cite{good1950probability}, when its inputs
and weights instantiate the constructive BP representation.

This is why the OR-gate language is useful: the FFN combines gathered
evidence into an updated conclusion. But the exact theorem is more
specific than the metaphor. The FFN implements the BP update exactly in
the sigmoid construction; production FFNs approximate or generalize this
role rather than literally instantiating the same binary rule in every
unit.

\subsection*{The Key-Value Memory Interpretation}

Geva et al.~\cite{geva2021ffn} show that FFN layers can be viewed as
key-value memories: the rows of $W_1$ are keys that detect input
patterns, and the columns of $W_2$ are values written to the residual
stream when those patterns activate. Each hidden unit detects a pattern
and contributes an update.

In the BP interpretation, these key-value memories resemble learned
factor potentials. The key determines when an update is relevant; the
value determines what direction is written back into the residual
stream. The full FFN layer integrates many such updates into the token's
next residual-stream state.

ROME~\cite{meng2022rome} provides suggestive evidence for this
interpretation: editing localized FFN parameters can change factual
associations in targeted ways. This does not by itself prove that FFN
weights are literal conditional-probability tables, but it is consistent
with the view that FFNs store semantically structured update directions.

\subsection*{Sigmoid vs.\ ReLU}

The formal proof requires sigmoid activation because $\updateBelief$ is
exactly $\sigma(\logit(m_0) + \logit(m_1))$: sigmoid is the inverse of
logit, and the composition gives the desired update. ReLU, GeLU, and
SwiGLU do not implement this binary log-odds algebra exactly.

Nevertheless, these activations may still support approximate or
generalized message-passing behavior. They preserve smooth, structured
transformations of residual-stream features, and empirical experiments
show that gradient descent can find approximately correct BP-like
solutions even with ReLU activations. In the \texttt{bayes-learner}
experiment, for example, the learned model reached validation MAE
$0.000752$~\cite{coppola2026transformers}.

Thus sigmoid gives the exact algebraic identity. Modern activations give
a practical approximation or extension, not the same formal guarantee.

\subsection*{Scaling}

For Llama~3 405B, there are
$L \times D_{ff} = 126 \times 53{,}248 \approx 6.7\mathrm{M}$
FFN hidden units across layers. The number of FFN evaluations per
forward pass is
$L \times D_{ff} \times W \approx 55\mathrm{B}$ at sequence length
$W = 8192$.

Within the BP interpretation, this large FFN budget reflects the cost of
implementing rich update rules over many superposed features. The
architecture devotes far more capacity to updating internal state than
to routing attention alone, which is consistent with the view that
inference is the dominant computational burden.